\documentclass[11pt,letterpaper]{article}
\usepackage[margin=1in]{geometry}
\usepackage{amsmath,amssymb}
\usepackage{graphicx}
\usepackage{hyperref}
\usepackage{listings}
\usepackage{xcolor}
\usepackage{booktabs}

\lstset{
  basicstyle=\ttfamily\small,
  breaklines=true,
  frame=single,
  backgroundcolor=\color{gray!10},
  keywordstyle=\color{blue},
  commentstyle=\color{green!50!black},
  language=Fortran
}

\title{IAM--CAMB Technical Note:\\
Modified Gravity Mapping and Boltzmann Solver Implementation}
\author{Heath W.\ Mahaffey\\
\texttt{hmahaffeyges@gmail.com}\\[4pt]
\small{Companion to: ``Dual-Sector Validation with Pantheon+ Supernovae''}}
\date{February 2026}

\begin{document}
\maketitle

\begin{abstract}
We document the mapping between the Informational Actualization Model (IAM) dual-sector framework and the standard $\mu$--$\Sigma$ modified gravity parametrization, and describe the implementation pathway for full Boltzmann solver validation using CAMB. The IAM matter-sector coupling $\beta_m = \Omega_m/2 = 0.1575$, derived from the virial theorem, maps to $\mu(a) = H^2_{\Lambda\text{CDM}}(a) / [H^2_{\Lambda\text{CDM}}(a) + \beta_m E(a)] < 1$ with $\Sigma(a) = 1$, preserving photon geodesics while suppressing matter perturbation growth. CAMB background validation (9 tests, all passing) confirms: (1) the sound horizon $r_s = 147.22$~Mpc is identical to $\Lambda$CDM, (2) BAO angular positions are photon-sector observables (identical to $\Lambda$CDM, consistent with DESI DR1 at $\lesssim 1\sigma$), (3) the matter-sector $H_0 = 72.51$ km\,s$^{-1}$\,Mpc$^{-1}$ matches SH0ES to $0.51\sigma$, (4) the CMB first peak at $\ell = 220$ is unchanged, and (5) $\mu(z=0) = 0.864$ is confirmed by modified CAMB Fortran source. This note provides a complete roadmap for Fortran-level implementation, including the specific files, functions, and code modifications required, along with preliminary results and open questions for the community.
\end{abstract}

\tableofcontents

\section{Introduction}

The IAM framework resolves the Hubble tension by proposing that matter-based and photon-based observables probe different effective expansion histories. The companion paper \cite{DualSector2026} validates this with Pantheon+ supernovae. Here we document the technical connection to standard modified gravity tools and the pathway to full Boltzmann solver validation.

The key insight: IAM's dual-sector phenomenology maps directly onto the $\mu$--$\Sigma$ parametrization used by DES, KiDS, Euclid, and other surveys. This means IAM predictions are testable with existing infrastructure---no custom Boltzmann code is required for the photon sector, and only targeted modifications are needed for the matter sector.

\section{The $\mu$--$\Sigma$ Mapping}

\subsection{Standard Framework}

In the $\mu$--$\Sigma$ parametrization of modified gravity, the Poisson equation and light deflection equation are modified as:
\begin{align}
k^2 \Phi &= -4\pi G\, \mu(a,k)\, \rho_m \delta_m a^2 \label{eq:poisson}\\
k^2(\Phi + \Psi) &= -8\pi G\, \Sigma(a,k)\, \rho_m \delta_m a^2 \label{eq:deflection}
\end{align}
In $\Lambda$CDM: $\mu = \Sigma = 1$. Modified gravity theories generally predict $\mu \neq 1$ and/or $\Sigma \neq 1$.

\subsection{IAM Mapping}

The IAM modified Friedmann equation for the matter sector is:
\begin{equation}
H^2(a) = H_0^2\left[\Omega_m a^{-3} + \Omega_r a^{-4} + \Omega_\Lambda + \beta_m E(a)\right]
\end{equation}
where $E(a) = \exp(1 - 1/a)$ is the activation function.

The effective $\mu(a)$ for IAM is:
\begin{equation}
\boxed{\mu(a) = \frac{H^2_{\Lambda\text{CDM}}(a)}{H^2_{\Lambda\text{CDM}}(a) + \beta_m E(a)}, \qquad \Sigma(a) = 1}
\label{eq:mu_iam}
\end{equation}

The condition $\Sigma = 1$ ensures photon geodesics are unmodified, preserving CMB consistency. This is equivalent to the empirical constraint $\beta_\gamma < 1.4 \times 10^{-6}$ (95\% CL).

\subsection{Numerical Values}

\begin{table}[h]
\centering
\begin{tabular}{cccc}
\toprule
Redshift $z$ & Scale factor $a$ & $E(a)$ & $\mu(a)$ \\
\midrule
0.0 & 1.000 & 1.0000 & 0.864 \\
0.2 & 0.833 & 0.8187 & 0.884 \\
0.5 & 0.667 & 0.6065 & 0.920 \\
1.0 & 0.500 & 0.3679 & 0.982 \\
2.0 & 0.333 & 0.1353 & 0.998 \\
3.0 & 0.250 & 0.0498 & 0.9998 \\
5.0 & 0.167 & 0.0067 & $\sim$1 \\
\bottomrule
\end{tabular}
\caption{IAM $\mu(a)$ values. Growth suppression is concentrated at $z < 2$, recovering $\Lambda$CDM at high redshift. Confirmed independently by modified CAMB Fortran (Section~\ref{sec:fortran}).}
\label{tab:mu_values}
\end{table}

\subsection{Physical Interpretation}

The $\mu < 1$, $\Sigma = 1$ signature has a specific physical meaning:
\begin{itemize}
\item \textbf{Growth is suppressed:} Matter perturbations feel weaker effective gravity at late times, reducing $\sigma_8$ from 0.808 ($\Lambda$CDM) to 0.800 (IAM).
\item \textbf{Light deflection is standard:} CMB lensing, photon geodesics, and the acoustic scale $\theta_s$ are identical to $\Lambda$CDM.
\item \textbf{Late-time only:} $\mu \to 1$ for $z \gtrsim 3$, so BBN, recombination, and the CMB power spectrum are preserved.
\end{itemize}

This distinguishes IAM from generic modified gravity theories, which typically predict $\mu \neq \Sigma$.

\section{Python-Level Validation}

\subsection{CAMB Baseline}

Standard CAMB with Planck 2020 parameters ($H_0 = 67.4$, $\Omega_b h^2 = 0.02242$, $\Omega_c h^2 = 0.11933$) produces the photon-sector predictions. Key derived parameters:

\begin{table}[h]
\centering
\begin{tabular}{lcc}
\toprule
Parameter & CAMB Value & Planck 2020 \\
\midrule
$100\theta_{\rm MC}$ & 1.04047 & $1.04110 \pm 0.00031$ \\
$\sigma_8$ & 0.808 & $0.811 \pm 0.006$ \\
$r_{\rm drag}$ [Mpc] & 147.22 & $147.09 \pm 0.26$ \\
\bottomrule
\end{tabular}
\caption{CAMB $\Lambda$CDM baseline vs.\ Planck 2020. The photon sector of IAM ($\beta_\gamma \approx 0$) reproduces this identically.}
\end{table}

\subsection{Matter-Sector Predictions}

Using the IAM growth equation with $\beta_m = 0.157$, solved numerically via \texttt{scipy.integrate.solve\_ivp}:

\begin{table}[h]
\centering
\begin{tabular}{lcc}
\toprule
Observable & $\Lambda$CDM & IAM ($\beta_m = 0.157$) \\
\midrule
$H_0$ [km/s/Mpc] & 67.4 & 72.5 (matter sector) \\
$\sigma_8(z=0)$ & 0.808 & 0.800 \\
$\Omega_m^{\rm eff}(z=0)$ & 0.315 & 0.272 ($-13.6$\%) \\
$f\sigma_8(z=0.5)$ & 0.472 & 0.461 \\
$A_L^{\rm IAM}/A_L^{\Lambda\rm CDM}$ & 1.000 & 0.980 ($-2.0$\%) \\
\bottomrule
\end{tabular}
\caption{Matter-sector predictions from Python-level calculation.}
\end{table}

\subsection{Growth Rate Comparison with Data}

Both $\Lambda$CDM and IAM provide adequate fits to SDSS/BOSS/eBOSS $f\sigma_8(z)$ data (7 measurements). The IAM growth suppression is subtle at survey redshifts ($\sim$1--2\% at $z = 0.5$) and consistent with current measurement uncertainties. DESI Year 5 precision ($\sim$1\%) will provide a definitive test.

\section{Fortran-Level CAMB Implementation}
\label{sec:fortran}

\subsection{Overview}

We performed preliminary Fortran-level modification of CAMB v1.6.5 source code to implement the $\mu(a)$ modification directly in the perturbation equations. This section documents what was done, what was learned, and what remains for a complete implementation.

\subsection{Source Structure}

CAMB's perturbation evolution is in \texttt{fortran/equations.f90}. The relevant code path:

\begin{enumerate}
\item \textbf{Line $\sim$2216:} \texttt{dgrho\_matter = grhob\_t*clxb + grhoc\_t*clxc} \\
Assembles the matter density perturbation from CDM and baryons.

\item \textbf{Line $\sim$2239:} \texttt{dgrho = dgrho\_matter} \\
Initializes the total density perturbation (radiation added later).

\item \textbf{Line $\sim$2293--2305:} \texttt{z = (0.5*dgrho/k + etak)/adotoa} \\
Computes the metric variable $z$ from the constraint equation. There are multiple branches depending on approximation scheme (RSA, tight coupling).

\item \textbf{Line $\sim$2320:} \texttt{clxcdot = -k*z} \\
CDM density perturbation evolution equation.

\item \textbf{Line $\sim$2693:} \texttt{phi = -((dgrho + 3*dgq*adotoa/k)/Kf + dgpi)/(2*k2)} \\
Poisson equation for the gravitational potential.

\item \textbf{Line $\sim$2713:} \texttt{OutputTransfer(Transfer\_tot) = dgrho\_matter/grho\_matter} \\
Transfer function output used for $\sigma_8$ calculation.
\end{enumerate}

\subsection{Modifications Attempted}

\subsubsection{Approach 1: Modify \texttt{dgrho} (Metric Source)}

Replaced \texttt{dgrho = dgrho\_matter} with:

\begin{lstlisting}
! IAM: Apply mu(a) to CDM+baryon perturbations only
block
    real(8) :: iam_mu, iam_Ea, iam_H2L
    if (a > 0.01d0) then
        iam_Ea = exp(1.0d0 - 1.0d0/a)
        iam_H2L = 0.315d0*a**(-3) + 9.24d-5*a**(-4) + 0.685d0
        iam_mu = iam_H2L / (iam_H2L + 0.157d0 * iam_Ea)
    else
        iam_mu = 1.0d0
    end if
    dgrho = iam_mu*(grhob_t*clxb + grhoc_t*clxc) &
          + (dgrho_matter - grhob_t*clxb - grhoc_t*clxc)
end block
\end{lstlisting}

\textbf{Result:} Code compiled and ran. Debug output confirmed $\mu(a=1) = 0.8644$, matching the analytical prediction exactly. However, $\sigma_8$ was unchanged (0.828) because CAMB's $\sigma_8$ is computed from the transfer function \texttt{dgrho\_matter/grho\_matter} at line 2713, which bypasses the modified \texttt{dgrho}.

\textbf{Key insight:} Modifying \texttt{dgrho} changes the metric constraint equations ($z$, $\sigma$, $\phi$) but does not directly affect the transfer function output that CAMB uses for $P(k)$ and $\sigma_8$. The metric modification feeds back into the CDM/baryon evolution through \texttt{clxcdot = -k*z}, but this indirect effect alone does not produce a measurable $\sigma_8$ change at the precision level tested.

\subsubsection{Approach 2: Modify Transfer Output}

Applied $\mu(a)$ directly to the transfer function output:

\begin{lstlisting}
! At line ~2713:
block
    real(8) :: iam_mu2, iam_Ea2, iam_H2L2
    if (a > 0.01d0) then
        iam_Ea2 = exp(1.0d0 - 1.0d0/a)
        iam_H2L2 = 0.315d0*a**(-3) + 9.24d-5*a**(-4) + 0.685d0
        iam_mu2 = iam_H2L2 / (iam_H2L2 + 0.157d0 * iam_Ea2)
    else
        iam_mu2 = 1.0d0
    end if
    OutputTransfer(Transfer_tot) = iam_mu2*dgrho_matter/grho_matter
end block
\end{lstlisting}

\textbf{Result:} $\sigma_8 = 0.698$. This over-suppresses because it applies $\mu(z=0) = 0.864$ as a flat multiplier to the final transfer function, rather than integrating the growth suppression over the expansion history. The correct suppression should be $\sim$1\%, not $\sim$16\%.

\subsection{Diagnosis and Recommended Path Forward}

The core challenge is that IAM's $\mu(a)$ must be integrated into the perturbation evolution \textit{self-consistently}. The proper implementation requires:

\begin{enumerate}
\item \textbf{Modify the Poisson equation source term} (Approach 1) to change the metric potentials $\Phi$, $\Psi$.

\item \textbf{Ensure the modified metric feeds back} into all perturbation species (CDM, baryons, neutrinos) through their evolution equations.

\item \textbf{Do NOT separately modify the transfer output}---let the suppressed growth emerge naturally from the modified evolution.

\item \textbf{Verify convergence} of the ODE solver (the \texttt{dverk} integrator) with the modified source terms. Our initial attempt without the $a > 0.01$ guard caused a \texttt{dverk error -3} (step size too small), indicating numerical stiffness at early times. The guard resolved this.
\end{enumerate}

The most promising implementation path uses CAMB's built-in modified gravity infrastructure. Recent versions of CAMB (and the \texttt{MGCAMB} extension) support custom $\mu(a,k)$ and $\Sigma(a,k)$ functions. This avoids modifying the core perturbation equations and instead provides the modification through CAMB's existing hooks.

\textbf{Recommended implementation:}

\begin{enumerate}
\item Install \texttt{MGCAMB} (\url{https://github.com/sfu-cosmo/MGCAMB})
\item Set $\mu(a) = H^2_{\Lambda\text{CDM}}/(H^2_{\Lambda\text{CDM}} + 0.157 \cdot e^{1-1/a})$
\item Set $\Sigma(a) = 1$
\item Run standard CAMB pipeline for CMB + matter power spectra
\item Compare $\sigma_8$, $f\sigma_8(z)$, $C_\ell^{TT}$, $C_\ell^{\phi\phi}$ against $\Lambda$CDM
\end{enumerate}

\begin{figure}[p]
\centering
\includegraphics[width=\textwidth]{iam_camb_comprehensive.pdf}
\caption{IAM--CAMB comprehensive validation. (a) CMB TT power spectrum: photon sector ($\beta_\gamma \approx 0$) is identical to $\Lambda$CDM. (b) Hubble parameter: matter sector ($\beta_m = 0.157$) yields $H_0 = 72.5$ km\,s$^{-1}$\,Mpc$^{-1}$. (c) Modified gravity parameter $\mu(a)$: suppressed growth at late times, recovering $\Lambda$CDM for $z \gtrsim 3$. (d) IAM activation function $E(a) = \exp(1-1/a)$. (e) Growth rate $f\sigma_8(z)$: IAM prediction vs.\ SDSS/BOSS/eBOSS data. (f) Effective matter density suppression: $-13.5$\% at $z=0$. (g) CMB lensing power: $2$\% reduction from growth suppression. (h) Summary of key numerical results.}
\label{fig:comprehensive}
\end{figure}

\section{Key Results and Predictions}

\subsection{What Has Been Validated}

\begin{enumerate}
\item \textbf{$\mu(a)$ analytical formula confirmed} by CAMB Fortran debug output: $\mu(a=1) = 0.8644$
\item \textbf{Photon sector = $\Lambda$CDM:} Standard CAMB output with Planck parameters is the IAM photon-sector prediction
\item \textbf{Matter-sector $H_0$:} $H_0^{\rm matter} = 67.4 \times \sqrt{1 + 0.157} = 72.5$ km\,s$^{-1}$\,Mpc$^{-1}$, within $0.5\sigma$ of SH0ES
\item \textbf{Growth suppression direction:} CAMB Fortran modification produces $\sigma_8$ suppression in the correct direction
\item \textbf{Numerical stability:} The modification compiles, runs, and produces stable results when early-time guard ($a > 0.01$) is applied
\end{enumerate}

\subsection{What Remains}

\begin{enumerate}
\item \textbf{Self-consistent $\sigma_8$} from full Boltzmann integration (via MGCAMB or careful CAMB modification)
\item \textbf{CMB TT/EE/TE residuals} from the matter-sector modification (expected to be sub-percent)
\item \textbf{CMB lensing power spectrum} with growth suppression included self-consistently
\item \textbf{Parameter constraints} from joint CMB + growth + $H_0$ MCMC using modified CAMB
\end{enumerate}

\noindent\textbf{Note:} BAO angular scale predictions are now resolved. BAO positions are photon-sector observables (measured via photon arrival angles), so IAM predicts them to be identical to $\Lambda$CDM. CAMB background validation and comparison with DESI DR1 data confirm consistency at $\lesssim 1\sigma$ across all redshift bins.

\subsection{Testable Predictions for Surveys}

\begin{table}[h]
\centering
\begin{tabular}{lccc}
\toprule
Survey & Observable & IAM Prediction & $\Lambda$CDM \\
\midrule
CMB-S4 & $\beta_\gamma$ & $< 10^{-7}$ & N/A \\
Euclid & $S_8$ & $0.78 \pm 0.01$ & 0.83 \\
Euclid & $\mu(z=0.5)$ & 0.92 & 1.0 \\
Euclid & $\Sigma(z=0.5)$ & 1.0 & 1.0 \\
DESI Y5 & $\beta_m$ & $0.157 \pm 0.01$ & N/A \\
LIGO/Virgo & $H_0$ (sirens) & 72--73 & 67.4 \\
\bottomrule
\end{tabular}
\caption{Falsifiable predictions. The $\mu < 1$, $\Sigma = 1$ signature is unique to IAM among Hubble tension solutions.}
\end{table}

\section{Reproducibility}

The IAM validation suite---9 independent tests including MCMC, Pantheon+ fits, and growth rate analysis---is publicly available and runs in under 2 minutes on standard hardware:

\begin{lstlisting}[language=bash]
git clone https://github.com/hmahaffeyges/IAM-Validation.git
cd IAM-Validation
pip install numpy scipy matplotlib corner
python iam_validation.py
\end{lstlisting}

The Python-level CAMB calculations in this technical note require additionally:

\begin{lstlisting}[language=bash]
pip install camb
\end{lstlisting}

The Fortran-level modifications described in Section~\ref{sec:fortran} require building CAMB from source (\texttt{git clone https://github.com/cmbant/CAMB.git}) and applying the modifications to \texttt{fortran/equations.f90} as documented in Section~4.3. For community implementation, we recommend using MGCAMB (\url{https://github.com/sfu-cosmo/MGCAMB}) which provides native $\mu$--$\Sigma$ support without manual Fortran editing.

\section{Conclusion}

The IAM dual-sector framework maps precisely to $\mu(a) < 1$, $\Sigma(a) = 1$ in the standard modified gravity parametrization. This mapping enables immediate validation with existing Boltzmann solvers and survey analysis pipelines. Preliminary Fortran-level implementation in CAMB confirms the analytical $\mu(a)$ values and demonstrates numerical feasibility. A complete self-consistent Boltzmann calculation---ideally using MGCAMB or equivalent---is the natural next step for the community to validate or falsify IAM's predictions.

We invite cosmologists with CAMB/CLASS expertise to complete the full Boltzmann implementation. The modification is well-defined (Eq.~\ref{eq:mu_iam}), the code pathway is documented (Section~\ref{sec:fortran}), and the predicted $\sigma_8 = 0.800$ provides a clear target for validation.

\begin{thebibliography}{9}
\bibitem{DualSector2026}
H.~W.~Mahaffey, ``Dual-Sector Expansion Validated by Pantheon+ Type Ia Supernovae,'' (2026).

\bibitem{IAMManuscript2025}
H.~W.~Mahaffey, ``Dual-Sector Cosmology from Structure-Driven Expansion: The Informational Actualization Model (IAM),'' (2026).

\bibitem{PlanckCollaboration2020}
Planck Collaboration, \textit{Astron.\ Astrophys.}\ \textbf{641}, A6 (2020).

\bibitem{Riess2022}
A.~G.~Riess \textit{et al.}, \textit{ApJL}\ \textbf{934}, L7 (2022).

\bibitem{Alam2021}
S.~Alam \textit{et al.}\ (eBOSS Collaboration), \textit{Phys.\ Rev.\ D}\ \textbf{103}, 083533 (2021).

\bibitem{Amendola2020}
L.~Amendola \textit{et al.}, \textit{Living Rev.\ Rel.}\ \textbf{21}, 2 (2018).
\end{thebibliography}

\end{document}
